% TOP_PROBLEM: {"id": 2100703, "title": "Open Problems in Integrable Systems — Integrability and Quantisation", "category_id": 11, "category": {"id": 11, "name": "dynamical_systems", "display_name": "Dynamical Systems", "description": "Problems about long-term behavior of deterministic systems, Hamiltonian dynamics, and periodic orbits.", "slug": "dynamical-systems", "order_index": 11, "created_at": "2026-07-31T15:26:25.671Z"}, "status": "partially_solved"}
% FIRST_POSTED: null
% ENTRY_KIND: statement_only
% Imported from: batch09.tex; UnsolvedMath ID 2100703
\documentclass[11pt]{article}
\usepackage[margin=25mm]{geometry}
\usepackage{fontspec}
\setmainfont{Latin Modern Roman}
\usepackage{amsmath,amssymb,mathtools}
\usepackage{unicode-math}
\setmathfont{Latin Modern Math}
\usepackage{xurl,hyperref}
\hypersetup{colorlinks=true,urlcolor=blue}
\setcounter{secnumdepth}{0}
\setlength{\parindent}{0pt}
\setlength{\parskip}{5pt}
\setlength{\emergencystretch}{3em}
\newcommand{\sref}[2]{\hyperlink{source:#1:#2}{[S#2]}}
\newcommand{\cataloguescope}[1]{\paragraph{Recorded target.}#1}
\begin{document}
% UnsolvedMath ID 2100703; source catalogue code AMR-020-0703
\setcounter{section}{2100702}
\section{Open Problems in Integrable Systems — Integrability and Quantisation}\label{2100703}
{\small\sffamily UnsolvedMath ID 2100703 · Dynamical Systems}
\subsection{Definitions and mathematical statement}

\paragraph{Shared notation.}$\mathbb N=\{1,2,\ldots\}$, and $\mathbb Z,\mathbb Q,\mathbb R,\mathbb C$ denote the integers, rationals, reals and complex numbers. Any other conventions are specified locally.


Let $(M^{2n},\omega)$, $n\geq2$ be a compact prequantizable Kähler manifold with a positive holomorphic line bundle $L$ of curvature $-i\omega$. For $k\to\infty$, put $\hbar=k^{-1}$ and $\mathcal H_k=H^0(M,L^{\otimes k})$. A Berezin--Toeplitz operator family has an expansion $P_k\sim\sum_{j\geq0}k^{-j}\Pi_k M_{p_j}\Pi_k$, where $\Pi_k$ is the orthogonal projection to $\mathcal H_k$ and $M_{p_j}$ multiplies by the smooth symbol $p_j$. Its principal symbol is $p_0$.\par Let $P_{1,k},\ldots,P_{n,k}$ be commuting self-adjoint Toeplitz families with principal-symbol map $F=(p_1,\ldots,p_n)$, where all $\{p_i,p_j\}=0$ and $dp_1,\ldots,dp_n$ are independent almost everywhere on $M$. A focus-focus block at a nondegenerate singular point is a four-dimensional transverse symplectic block whose commuting quadratic normal forms, after an invertible linear change of components, are
\[q_1=x_1\xi_1+x_2\xi_2,\qquad q_2=x_1\xi_2-x_2\xi_1.\]
Here $(x_1,x_2,\xi_1,\xi_2)$ are Darboux coordinates on that block. In dimension four these are the rank-zero focus-focus model; in higher dimensions they occur alongside the remaining regular directions and nondegenerate transverse blocks. Focus-focus singularities in the question include these higher-dimensional cases. The joint spectrum consists of $n$-tuples of simultaneous eigenvalues. An error $O(k^{-\infty})$ means an error $O(k^{-A})$ for every $A>0$.
\paragraph{Statement recorded in the supplied source.}
Construct complete Bohr--Sommerfeld quantization rules near focus-focus singular fibers in Berezin--Toeplitz quantization: geometric integrality conditions, including corrections to every order in $k^{-1}$, that are necessary and sufficient for $E\in\mathbb R^n$ to belong to the joint spectrum modulo $O(k^{-\infty})$. \sref{2100703}{2}
\subsection{Short English statement}
Give all-order geometric quantization conditions describing the joint spectrum near a focus-focus singularity for commuting Toeplitz operators.
\subsection{Sources}
\begin{description}
\item[\hypertarget{source:2100703:1}{S1}] ULAM.AI and the UnsolvedMath Contributors, UnsolvedMath: A Curated Collection of Open Mathematics Problems, record 2100703 (AMR-020-0703). CC BY 4.0. \url{https://huggingface.co/datasets/ulamai/UnsolvedMath}.
\item[\hypertarget{source:2100703:2}{S2}] Alexey Bolsinov, Vladimir S. Matveev, Eva Miranda and Serge Tabachnikov, Open Problems, Questions, and Challenges in Finite-Dimensional Integrable Systems (2018), Problem 7.3 and its complete preceding definitions. Original question and full discussion. \url{https://arxiv.org/abs/1804.03737}.

\end{description}
\label{end:2100703}
\end{document}
